%% Environmental Modelling & Software -- elsarticle, author-year (Harvard)
%% Supplementary Material: worked examples for the non-parametric tests used
%% in the main text (Section 3.5, Validation).
\documentclass[12pt,authoryear]{elsarticle}

\usepackage[T1]{fontenc}
\usepackage[utf8]{inputenc}
\usepackage{lmodern}
\usepackage{amsmath,amssymb}
\usepackage{graphicx}
\usepackage{booktabs}
\usepackage[hidelinks]{hyperref}
\usepackage{microtype}
\usepackage[margin=2.5cm]{geometry}
\biboptions{authoryear,round,semicolon}

\begin{document}

\begin{frontmatter}
\title{Supplementary Material: Worked Examples for the Statistical Tests
       Used in the Main Text}
\author[unipa]{Cicero Martins Jr.\corref{cor1}}
\ead{cicero.martinsjr@unipa.it}
\author[unipa]{Fulvio Capodici}
\author[unipa]{Mauro De Marchis}
\author[unipa]{Giuseppe Ciraolo}
\cortext[cor1]{Corresponding author}
\affiliation[unipa]{organization={Dipartimento di Ingegneria,
                    Universit\`a degli Studi di Palermo},
                    city={Palermo}, country={Italy}}
\end{frontmatter}

\noindent This document accompanies
\emph{Shoreline compactness predicts Sentinel-1 water-area monitoring
reliability across reservoirs worldwide} and gives a self-contained,
hand-checkable worked example for each of the four quantitative tools
introduced in Section~3.5: the accuracy metric (KGE) and the three
non-parametric tests applied to it. Each example uses a small, real subset of
the reservoirs analysed in the main text so the mechanics of the calculation
can be followed by hand, then compared against a statistical package. Because
each subset is far smaller than the full sample used in the paper ($n=62$ or
$n=59$, depending on the test), the resulting $p$-values here are not
significant. That is expected and is not a discrepancy: these examples
illustrate \emph{how} each statistic is computed, not a replication of the
main text's headline results.

\section*{S1. Kling-Gupta Efficiency (KGE) \citep{Gupta2009}}

\textbf{Used for:} the core accuracy score throughout the paper, comparing
a SAR-derived monthly area series against a reference (JRC or PlanetScope).
It is the quantity that Sections~S2--S4 below then compare or correlate.

\textbf{Why this metric:} the two obvious alternatives are the Pearson
correlation $r$ alone, and the older Nash-Sutcliffe Efficiency (NSE). Both
conflate several distinct ways a simulated series can fail to match an
observed one. A series can track the timing of every rise and fall
perfectly ($r=1$) while still being consistently too high, or while swinging
much more sharply than the truth, and a single correlation or NSE value
cannot tell these failure modes apart. KGE decomposes accuracy into three
independent components so each can be inspected on its own.

\textbf{Calculation.} For observed series $o$ and simulated series $s$ of
equal length,
\[
\mathrm{KGE} = 1 - \sqrt{(r-1)^2 + (\alpha-1)^2 + (\beta-1)^2},
\]
where $r$ is the Pearson correlation between $o$ and $s$ (captures timing
and shape), $\alpha = \sigma_s/\sigma_o$ is the variability ratio (captures
whether $s$ swings by the right amount), and $\beta = \mu_s/\mu_o$ is the
bias ratio (captures whether $s$ is centred correctly). All three equal 1
under perfect agreement, so $\mathrm{KGE}=1$ is a perfect score, and it
decreases as any one component departs from 1, regardless of the other two.

\textbf{Worked example} (Cecita, 7 consecutive months of JRC-referenced
monthly area, ha):

\begin{center}
\begin{tabular}{lrrrrrrr}
\toprule
& 2015-03 & 04 & 05 & 06 & 07 & 08 & 09 \\
\midrule
JRC (obs.), $o$    & 773.5 & 780.2 & 768.8 & 768.4 & 739.4 & 684.4 & 544.9 \\
VV Otsu (sim.), $s$ & 701.3 & 878.4 & 879.7 & 791.6 & 712.8 & 652.6 & 569.5 \\
\bottomrule
\end{tabular}
\end{center}

$\bar{o}=722.8$, $\bar{s}=740.8$, $\sigma_o=78.8$, $\sigma_s=107.1$, and the
Pearson correlation of the two rows is $r=0.812$. So
\[
\alpha = \frac{107.1}{78.8} = 1.359, \qquad
\beta = \frac{740.8}{722.8} = 1.025,
\]
\[
\mathrm{KGE} = 1 - \sqrt{(0.812-1)^2 + (1.359-1)^2 + (1.025-1)^2}
             = 1 - \sqrt{0.0353 + 0.1289 + 0.0006} = 0.594.
\]
The correlation alone looks reasonable ($r=0.81$, the seasonal drawdown is
captured), but $\alpha=1.36$ shows VV Otsu over-swings relative to JRC by
about 36\% in this window, which is what actually drags the composite score
down to 0.59. This is the diagnostic value of the decomposition: a
correlation-only or NSE score would not have distinguished a timing problem
from an amplitude problem. \texttt{scipy}, applied to the same seven-month
window, reproduces $r$, $\alpha$, and $\beta$ to the same precision.

\subsection*{A note on $p$-values and significance}

Sections~S2--S4 below each return two numbers: a statistic ($\rho$, $W$, or
$H$) that summarises the pattern found in the sample, and a $p$-value that
judges how surprising that pattern would be under a specific no-effect
scenario, the \emph{null hypothesis} ($H_0$): no monotonic relationship
(Spearman), no systematic paired difference (Wilcoxon), or no difference
between groups (Kruskal-Wallis). Formally, the $p$-value is the probability
of obtaining a statistic at least as extreme as the one observed if $H_0$
were exactly true. A small $p$-value means the observed pattern would be
unlikely to arise by chance alone, and is treated as evidence against
$H_0$. Throughout the paper we use the conventional threshold $p<0.05$
\citep{Fisher1925}: results below it are called \emph{statistically
significant} (the pattern is unlikely to be a sampling artefact), and
results at or above it are called \emph{not significant} (NS), meaning the
data do not rule out $H_0$, not that no effect exists. This last point
matters for the worked examples below: each uses only 7--8 reservoirs to
keep the arithmetic tractable by hand, and at that sample size even a real,
large effect will usually not clear $p<0.05$. That is why every
worked-example $p$-value in Sections~S2--S4 is NS while the corresponding
full-sample result in the paper is significant, the effect is the same,
only the statistical power to detect it differs. KGE (Section~S1) is not a
hypothesis test and has no associated $p$-value: it is a descriptive
accuracy score bounded above by 1, not a statistic being compared to a null
distribution.

\section*{S2. Spearman rank correlation \citep{Spearman1904}}

\textbf{Used for:} testing whether the area-to-perimeter ratio (A/P)
predicts surface-area skill (KGE), Section~4.1, and similar monotonic
relationships throughout the paper.

\textbf{Why this test:} Spearman's $\rho$ tests for a monotonic relationship
without assuming it is linear or that either variable is normally
distributed. This matters here because A/P is expected to set a
\emph{ceiling} on KGE (rising sharply then saturating), not a straight-line
trend, so a rank-based measure is the appropriate one.

\textbf{Calculation.} Both variables are converted to ranks (1 = smallest),
and the ordinary correlation coefficient is computed on the ranks rather
than the raw values:
\[
\rho = 1 - \frac{6 \sum_i d_i^2}{n(n^2-1)}, \qquad
d_i = \mathrm{rank}(X_i) - \mathrm{rank}(Y_i).
\]

\textbf{Worked example} (7 reservoirs spanning the A/P range):

\begin{center}
\begin{tabular}{lrrrrr}
\toprule
Reservoir & A/P (m) & KGE & rank(A/P) & rank(KGE) & $d^2$ \\
\midrule
Long\_Lake       & 62  & 0.251 & 1 & 1 & 0 \\
Sau              & 84  & 0.878 & 2 & 5 & 9 \\
Panneciere       & 111 & 0.743 & 3 & 4 & 1 \\
Cecita           & 125 & 0.953 & 4 & 6 & 4 \\
Chivor           & 158 & 0.693 & 5 & 2 & 9 \\
Baisha           & 209 & 0.711 & 6 & 3 & 9 \\
Hassan\_Addakhil & 267 & 0.963 & 7 & 7 & 0 \\
\bottomrule
\end{tabular}
\end{center}

With $n=7$ and $\sum d_i^2 = 32$:
\[
\rho = 1 - \frac{6 \times 32}{7(49-1)} = 1 - \frac{192}{336} = 0.429.
\]
For reference, \texttt{scipy.stats.spearmanr} on the same seven pairs
returns $\rho = 0.429$, $p = 0.34$ (not significant at $n=7$, as expected).
On the full 62-reservoir set used in the paper, the same calculation gives
$\rho = 0.51$, $p = 3\times10^{-5}$ (Section~4.1).

\section*{S3. Wilcoxon signed-rank test \citep{Wilcoxon1945}}

\textbf{Used for:} the paired, per-reservoir comparison between the
per-scene SVM and VV Otsu detectors, Section~4.2.

\textbf{Why this test:} it is the non-parametric analogue of the paired
$t$-test, appropriate because the paired KGE differences
($\mathrm{KGE}_{\text{SVM}} - \mathrm{KGE}_{\text{VV}}$) are strongly
right-skewed (a few very large SVM wins alongside many smaller VV wins)
rather than normally distributed, which would violate the assumption behind
a paired $t$-test.

\textbf{Calculation.} For each pair, compute the signed difference
$d_i = \mathrm{KGE}_{\mathrm{SVM},i} - \mathrm{KGE}_{\mathrm{VV},i}$, rank the
\emph{absolute} differences $|d_i|$ from smallest to largest, then reattach
each rank's original sign. Sum the ranks of the positive differences ($W^+$)
and of the negative differences ($W^-$). The test statistic is
$W = \min(W^+, W^-)$.

\textbf{Worked example} (7 reservoirs, ordered by $|d_i|$):

\begin{center}
\begin{tabular}{lrrrrr}
\toprule
Reservoir & KGE$_{\text{VV}}$ & KGE$_{\text{SVM}}$ & $d_i$ & $|d_i|$ rank & signed rank \\
\midrule
Cardinia       & 0.855  & 0.874 & $+0.019$ & 1 & $+1$ \\
Karapuzha      & 0.773  & 0.723 & $-0.050$ & 2 & $-2$ \\
Hubbard\_Creek & 0.755  & 0.504 & $-0.251$ & 3 & $-3$ \\
Saint\_Cassien & 0.678  & 0.347 & $-0.331$ & 4 & $-4$ \\
Sau            & 0.515  & 0.878 & $+0.363$ & 5 & $+5$ \\
Ambuklao       & 0.048  & 0.859 & $+0.811$ & 6 & $+6$ \\
Baisha         & $-0.320$ & 0.711 & $+1.031$ & 7 & $+7$ \\
\bottomrule
\end{tabular}
\end{center}

$W^+ = 1+5+6+7 = 19$, $W^- = 2+3+4 = 9$, so $W = \min(19,9) = 9$.
\texttt{scipy.stats.wilcoxon} on the same seven pairs returns the same
statistic and $p = 0.47$ (not significant at $n=7$). Note that four of the
seven reservoirs favour VV (negative $d_i$) while only three favour SVM, yet
the rank sums are closer together than the raw count suggests: the SVM's
three wins carry the three \emph{largest} ranks (5, 6, 7), which is exactly
the asymmetric-magnitude pattern discussed in Section~4.2. On the full
59-reservoir paired set used in the paper, the same calculation gives
$p = 0.37$.

\section*{S4. Kruskal-Wallis test \citep{KruskalWallis1952}}

\textbf{Used for:} testing whether the KGE residual, after accounting for
A/P, still differs across the four climate-zone families, Section~4.5.

\textbf{Why this test:} it is the non-parametric analogue of one-way ANOVA
for comparing more than two independent groups, appropriate given the small
and uneven group sizes ($n=12$--19 per biome in the full analysis) where
normality and equal variance within each group cannot be assumed.

\textbf{Calculation.} First isolate the climate signal from the geometric
one: fit a linear regression of KGE on A/P across all reservoirs, and take
the residual $e_i = \mathrm{KGE}_i - (\hat a + \hat b \cdot \mathrm{A/P}_i)$
for each reservoir. Pool the residuals from all groups, rank them 1 to $N$,
and sum the ranks within each group ($R_j$). The test statistic is
\[
H = \frac{12}{N(N+1)} \sum_j \frac{R_j^2}{n_j} - 3(N+1),
\]
which is approximately $\chi^2$-distributed with $k-1$ degrees of freedom
($k=4$ biome families here, so 3 d.f.).

\textbf{Worked example} (2 reservoirs per biome; the fitted line on this
8-point subset is $\mathrm{KGE} = 0.499 + 0.00108\times\mathrm{A/P}$):

\begin{center}
\begin{tabular}{llrrr}
\toprule
Reservoir & Biome & A/P (m) & KGE & residual $e_i$ \\
\midrule
Boadella    & Mediterranean          & 93  & 0.832 & $+0.232$ \\
Bilancino   & Mediterranean          & 160 & 0.683 & $+0.011$ \\
Amir\_Kabir & Semi-arid/arid         & 92  & 0.322 & $-0.276$ \\
Recoleta    & Semi-arid/arid         & 237 & 0.933 & $+0.177$ \\
Conestogo   & Temperate/continental  & 123 & 0.729 & $+0.097$ \\
Cardinia    & Temperate/continental  & 156 & 0.874 & $+0.206$ \\
Ambuklao    & (Sub)tropical          & 118 & 0.859 & $+0.232$ \\
Asolamandha & (Sub)tropical          & 255 & 0.938 & $+0.163$ \\
\bottomrule
\end{tabular}
\end{center}

Ranking the 8 residuals and summing within each 2-member group gives
$H = 1.17$ on 3 d.f., $p = 0.76$ (not significant, as expected for
$n=8$ split four ways). \texttt{scipy.stats.kruskal} on the same eight
residuals returns the same value. On the full 62-reservoir set used in the
paper, the same procedure (pooling all reservoirs' A/P-regression residuals
and grouping by biome) gives $H$ corresponding to $p = 0.012$
(Section~4.5), which is significant: unlike this small illustrative subset,
the full sample shows that climate carries a genuine signal beyond what A/P
alone explains.

\bibliographystyle{elsarticle-harv}
\bibliography{references}

\end{document}
